\documentclass[12pt,fleqn]{article}
\usepackage[a4paper, compat2]{geometry}
\usepackage[dvips]{epsfig}
\usepackage{amsfonts}
\usepackage{amssymb}
\usepackage{amsmath}
\usepackage{latexsym}
\usepackage{graphicx}
\usepackage[authoryear]{natbib}
\usepackage{setspace}
\usepackage{rotating}
\usepackage{multirow} % per mettere pi� righe nella stessa casella di una tabella
\usepackage{verbatim}
\usepackage{authblk}
\usepackage[english]{babel}
\usepackage[T1]{fontenc}
\usepackage[latin1]{inputenc}
\usepackage{subfigure}
\usepackage{pgf,pgfarrows,pgfnodes,pgfautomata,pgfheaps}

%\usepackage[dvipdf]{graphicx}


\newcommand{\mbf}[1]{\mathbf{#1}} %grassetto nelle formule
\newcommand{\ud}{\mathrm{d}} %differenziale
\newcommand{\bb}{\mathbb{R}} %numeri reali
\newcommand{\ra}{\rightarrow}
\newcommand{\beq}{\begin{equation}}
\newcommand{\eeq}{\end{equation}}
\newcommand{\bea}{\begin{eqnarray}}
\newcommand{\eea}{\end{eqnarray}}
\newcommand{\ba}{\begin{array}}
\newcommand{\ea}{\end{array}}
\newcommand{\bi}{\begin{itemize}}
\newcommand{\ei}{\end{itemize}}
\newcommand{\ben}{\begin{enumerate}}
\newcommand{\een}{\end{enumerate}}
\newcommand{\bfl}{\begin{flalign}}
\newcommand{\efl}{\end{flalign}}
\newcommand{\nn}{\nonumber}
\newcommand{\fn}[1]{\footnote{#1}}
\renewcommand{\r}{\right}
\renewcommand{\l}{\left}
\long\def\symbolfootnote[#1]#2{\begingroup\def\thefootnote{\fnsymbol{footnote}}\footnote[#1]{#2}\endgroup}

\newtheorem{es}{Ex.}

\title{Training Exercises}
\date{}
\author{}

%if no date, puts the date of today

\begin{document}
\maketitle
\section*{Sequences}
\begin{es}
Compute $\lim_{n\ra\infty} a_n$ for the following sequences
\beq	\nn
a_n=\sqrt[n]{n}\qquad\qquad a_n=\frac{n}{c^n}\quad c>1\qquad\qquad a_n=\frac{n}{n+1}\qquad\qquad\qquad a_n = x^n \quad x\in\mathbb R
\eeq
\beq	\nn
a_n=\sqrt[n]{2^n+3^n}\qquad\qquad a_n=\frac{\sin n}{\sqrt{n}}\qquad\qquad a_n=\sqrt[n]{n^{\alpha}}\quad \alpha>0\qquad\qquad a_n=\sqrt{n+1}-\sqrt n
\eeq
\end{es}
\begin{es}
Consider the sequence defined as 
\beq\nn
\l\{\ba{ccc}
a_1&=&0\\
&&\\
a_{k+1}&=&\sqrt{2+a_k}\quad \mbox{for}\quad k>1
\ea\r.
\eeq
Compute $\lim_{k\ra\infty} a_k$.\\
\end{es}
\section*{Limits}
\begin{es}
Compute the following limits
\beq\nn
\lim_{x\ra 0^+}\frac {2}{1+e^{-1/x}}\hskip 2cm \lim_{x\ra +\infty}\frac{2x+1}{x-1}\hskip 2cm \lim_{x\ra 1}\frac{x+1-2\sqrt x}{(x-1)^2}
\eeq

\beq\nn
\lim_{x\ra +\infty}\frac{2x^2-5+\sqrt{x^4-3x-1}}{x-1+\sqrt{x^4+x-2}}\hskip 3.4cm \lim_{x\ra +\infty}\sqrt{x^2+9x+3}-x
\eeq

\beq\nn
\lim_{x\ra 3^+}\frac{x-3}{\sqrt x-\sqrt 3}\hskip 2cm \lim_{x\ra +\infty}\frac{\log(1+2e^x)}{\sqrt{1+x^2}}\hskip 1cm \lim_{x\ra +\infty}\frac{x+4}{x^2+x+5}
\eeq

\beq\nn
\lim_{x\ra +\infty}\sqrt{x^2+4}-x\hskip 1.5cm\lim_{x\ra 0} x^2 \cos\frac 1 x\hskip 2cm\lim_{x\ra 0} x^2 \sin\frac 1 x\hskip 1.5cm\lim_{x\ra 3}\frac{x^2-9}{x-3} 
\eeq

\beq\nn
\lim_{x\ra 2}\frac{x-2}{\sqrt{2x^2+1}-3}\hskip 2cm\lim_{x\ra 0}\frac{\sin(2x)\sin(3x)}{x^2}
\eeq
\end{es}
\begin{es}
Compute, if they exist, the following limits
\beq\nn
\lim_{x\ra +\infty}\frac{\sqrt{9x^2+4x+5}}{x}\hskip 2cm \lim_{x\ra -\infty}\frac{\sqrt{9x^2+4x+5}}{x}
\eeq

\beq\nn
\lim_{x\ra 0}\frac{\sqrt{5x^2+x^4}}{x}\hskip 3cm \lim_{x\ra 0}\frac{|\sin x|}{x}\hskip 2cm 
\eeq

\beq\nn
\lim_{x\ra 0^+}\frac{\sqrt{x^2+x^3}}{x}\hskip 3cm \lim_{x\ra 0^-}\frac{\sqrt{x^2+x^3}}{x}\hskip 2cm \lim_{x\ra 0}\frac{\sqrt{x^2+x^3}}{x}
\eeq
\end{es}
\begin{es}
Compute the following limits, using de L'H\^opital's rule
\beq\nn
\lim_{x\ra 0}\frac{x-\sin x}{x^3}\hskip 3.5cm \lim_{x\ra 0}\frac{\tan x-x}{x^3}
\eeq

\beq\nn
\lim_{x\ra 1}\frac{x^3-3x+2}{x^3+x^2-5x+3}\hskip 2cm \lim_{x\ra -2}\frac{x^2+3x+2}{2x^2-8}
\eeq
\end{es}
\begin{es}
Compute the following limits
\beq\nn
\lim_{x\ra 1}\frac{\sqrt{1+3x}-2}{x-1}\hskip 2cm \lim_{x\ra -1}\frac{2x^2+7x+5}{3x^2+5x+2}
\eeq

\beq\nn
\lim_{x\ra 0}\frac{1-\cos (4x)}{x^2}\hskip 2cm \lim_{x\ra 2}\frac{x^2-4}{x-2}\hskip 2cm \lim_{x\ra 1}\frac{x^{1/3}-1}{x-1}
\eeq

\beq\nn
\lim_{x\ra +\infty}\frac{\cos(3x)}{x}\hskip 2cm \lim_{x\ra 0}x^2\l(1+\cos\frac 1 x \r)\hskip 2cm \lim_{x\ra +\infty}\frac{\cos x^2}{x}
\eeq

\beq\nn
\lim_{x\ra 0}\sqrt x\sin x\cos x\hskip 2cm \lim_{x\ra 0}\frac{\sin(-5x)}{7x}\hskip 2cm \lim_{x\ra 0}\frac{\cos x -1}{\sin^2 x}
\eeq
\end{es}
\newpage
\section*{Functions in $\mathbb {R}$}
\begin{es}
Determine the domain of the following functions
\beq\nn
f(x)=(x^3+1)^{-1}\hskip 2cm f(x)=\sqrt{1-\frac 1 x}\hskip 2cm f(x)=\sqrt{5-x}+\frac{1}{\sqrt{x+1}}
\eeq

\beq\nn
f(x)=\sin x+\tan x\hskip 2cm f(x)=\log(1-x^2)\hskip 2cm f(x)=\sqrt{4-x^2}
\eeq
\end{es}
\begin{es}
For the following functions study the domain, limits at the extremes of the domain, asymptots, local and global extrema, concavity.
Draw a graph of the function.
\beq
f(x)=4x-x^2+5
\eeq

\beq
f(x)=(2+\cos^2 x) \sin x
\eeq

\beq
f(x)=xe^{-x}
\eeq

\beq
f(x)=x+\frac 1 x
\eeq

\beq
f(x)=x-\log x
\eeq

\beq
f(x)=3x^2-4x-5
\eeq

\beq
f(x)=x-1+\sin x
\eeq

\beq
f(x)=\frac{1}{x^2+4}-3x+4
\eeq

\beq
f(x)=\log\l(\frac{x^{1/3}}{3x-1}\r)
\eeq

\beq
f(x) = 2x+3(e^x-2)^{2/3}
\eeq

\beq
f(x) = \l|x-\frac{x}{\log x}\r|
\eeq

\beq
f(x) = 1-(\log x)^3
\eeq

\beq
f(x)=e^{1/x}+\frac{\sin x}{x}
\eeq

\beq
f(x)=xe^{\frac{x-1}{x+3}}
\eeq

\beq
f(x)=\sqrt{x^2+1}-\sqrt x
\eeq

\beq
\log (|x+1|+e^x)
\eeq

\beq
f(x)=x^5-10x^3+25x
\eeq

\beq
f(x)=\frac{x+1}{x^2-9}
\eeq

\beq
f(x)=\frac{x^2+1}{x^2-1}
\eeq

\beq
f(x)=\frac{x^2+1}{(x-2)(x-4)}
\eeq

\beq
f(x)=\frac{x}{x^2+1}
\eeq

\beq
f(x)=\frac{x^2-3}{x-2}
\eeq

\beq
f(x)=\frac{x^2+3x-3}{x-1}
\eeq

\beq
f(x)=x^2-7x+6
\eeq

\beq
f(x)=x^3-x
\eeq

\beq
f(x)=\cos(2x)
\eeq

\beq
f(x)=\log|x|
\eeq

\beq
f(x)=\frac{x-1}{x+2}
\eeq

\beq
f(x)=\sqrt{x^2+1}
\eeq

\beq
f(x)=\frac{1}{x^2+1}
\eeq
\end{es}
\begin{es}
Determine $a$, $b$, $c$ and $d$ such that the function $f(x)=ax^3+bx^2+cx+d$ passes through the points $(0,3)$ and $(2,5)$ and has stationary points in $x=1/3$ and $x=1$.
\end{es}
\end{document}
